\documentclass[11pt,a4paper]{letter}
\usepackage[T1]{fontenc}
\usepackage{mathptmx}
\usepackage[margin=2.5cm]{geometry}
\usepackage{xcolor}
\usepackage[hidelinks]{hyperref}
\newcommand{\TBD}[1]{\textcolor{red}{\textbf{[#1]}}}

\signature{Keisuke Nishioka\\
Institute of Continuum Mechanics, Leibniz Universit\"at Hannover\\
\href{mailto:keisuke.nishioka@stud.uni-hannover.de}{keisuke.nishioka@stud.uni-hannover.de}\\
ORCID: \href{https://orcid.org/0009-0001-9360-1336}{0009-0001-9360-1336}}
\address{Institute of Continuum Mechanics\\Leibniz Universit\"at Hannover\\An der Universit\"at 1\\30823 Garbsen, Germany}
\date{\TBD{date}}

\begin{document}
\begin{letter}{Guest Editors, Special Issue ``Artificial Intelligence and Computational Mechanics for Materials Design''\\
Extreme Mechanics Letters}

\opening{Dear Dr.~Tac, Dr.~Moreno-Mateos, Dr.~Hossain and Prof.~Zhang,}

On behalf of all co-authors, I submit the manuscript
``\TBD{title}''
by K.~Nishioka, F.~Klempt, H.~Geisler, R.~Mukherjee, N.~Heine, K.~Doll-Nikutta, M.~Stiesch, S.~P.~Szafra\'{n}ski, M.~Soleimani and P.~Junker
for the special issue ``Artificial Intelligence and Computational Mechanics for Materials Design'' (VSI: AI \& CM), following the invitation sent to Prof.~Junker.

Biofilms are living materials whose composition, and with it their mechanical behaviour, is set by interactions between species that cannot be measured directly.
We identify these interactions as a physics-constrained inverse characterization problem: a thermodynamically consistent multiphase model, derived from the extended Hamilton principle, fixes the structure, and Bayesian inference determines its 15 interaction coefficients from in-vitro time courses of a five-species oral biofilm.
The forward model is compiled in JAX and evaluated for all particles in parallel on a GPU, which makes it affordable to repeat every inference stage with independent seeds and to accept a stage only after explicit checks of mixing, seed agreement and posterior mass at the prior bounds.
\TBD{Two sentences on the main result: which interactions the data determine, and the in-silico knockout prediction.}

We believe the work fits the special issue because it combines a mechanistic, variationally derived continuum model with data-driven probabilistic inference for the inverse characterization of a soft biological material, and because it makes explicit which conclusions the data support---the kind of physics-based guidance that the call identifies as essential for data-driven methods in materials science.

The manuscript has not been published and is not under consideration elsewhere.
All authors have approved the submission \TBD{confirm}.
Code, run configurations, convergence records and posterior samples are archived at Zenodo (\TBD{DOI}).
The authors declare no competing interests.

\closing{Sincerely,}

\end{letter}
\end{document}
